\subsection{Model and its basic characteristics}\label{subsec:r-model}

The model achieved 90.8\% accuracy on the training set and 87.6\% accuracy
on the validation set. Grad-CAM reveals that the network's attention is focused
on birds, effectively ignoring complex backgrounds and occlusions. The attention
maps cover the entire bird, indicating that the extracted phenotypic vectors
can be regarded as representations of the shape, plumage, and colour of species (\cref{fig:gradcam}).

\begin{figure}[htbp]
    \centering
    \includegraphics[width=\textwidth]{gradcam}
    \caption{Grad-CAM reveals that the network's attention is consistently
    focused on birds, ignoring backgrounds and occlusions.}
    \label{fig:gradcam}
\end{figure}

The global mean pairwise angle across all birds is approximately
1.57066 radians ($\approx \frac{\pi}{2}$), with a low variance
($\approx 0.00968$). In contrast, the intra-taxa distributions (families and orders) exhibit
significantly lower mean angles ($\approx1.2$) but higher variances
(peaking at $\approx0.02\textminus0.03$) (\cref{fig:angles}).

A strong positive correlation exists between taxa size
and mean pairwise angle (Orders: $\rho=0.91$; Families: $\rho=0.75$).
Mean pairwise angle shows weak correlation with angle variance
(Orders: $\rho=0.33$; Families: $\rho=0.47$).
Conversely, taxa size shows weak or no correlation with angle variance
(Orders: $\rho=0.24,p>0.05$; Families: $\rho=0.34$).

\begin{figure}[htbp]
    \centering
    \includegraphics[width=\textwidth]{angles_analysis}
    \caption{(a-f) The relationship between taxa size and mean pairwise angle,
        taxa size and pairwise angle variance, as well as mean pairwise angle and pairwise
        angle variance at both order and family levels. Spearman's rank coefficients ($\rho$)
        and p-values are listed for each figure; (g-h) the distribution of mean pairwise angle
        and pairwise angle variance of families and orders, the mean values of all birds are
        marked as the red dash line.}
    \label{fig:angles}
\end{figure}

\subsection{Similarity clustering}\label{subsec:r-similarity-clustering}

The clustering process results in a comprehensive hierarchical clustering output.
The agglomerative clustering technique applied to the cosine similarity
measures of the weight vectors yielded a dendrogram that illustrates the
relationships between the different avian species based on their
morphological features learned by the ResNet34 model.

In the taxonomic consistence analysis, we identified a total of 391
branches with high taxonomic consistency at the family level and 94
branches at the order level. Additionally, we found 474 family-level
outlier species and 533 order-level outlier species (\cref{fig:cluster}).

\begin{figure}[htbp]
  \centering
  \begin{minipage}{0.45\textwidth}
    \centering
    \includegraphics[width=1.05\textwidth]{cluster-1}
  \end{minipage}
  \hfill
  \begin{minipage}{0.45\textwidth}
    \centering
    \includegraphics[width=\textwidth]{cluster-2}
  \end{minipage}
  \caption{The unrooted clustering result with high-purity branches collapsed.}
\label{fig:cluster}
\end{figure}

\subsection{Disparity analysis}\label{subsec:r-disparity-analysis}

Regarding the relationship between order-level disparity and diversity, the
Spearman's coefficient is at 0.966, with a p-value of $1.45\times10^{-24}$.
For families, the Spearman's coefficient was 0.908,
with a p-value of $3.56\times10^{-83}$.

The stretched exponential model receives the lowest AIC (-251.33) at the order
level, followed by the Hill equation (-250.02), without significant difference
($\Delta AIC = 1.31$). The models are respectively $f(x) = 1 - \exp(-0.1420(x - 1)^{0.3081})$
and $f(x) = \frac{(x - 1)^{0.3973}}{7.4957 + (x - 1)^{0.3973}}$.
Orders with the highest disparity residuals are Falconiformes, Cuculiformes,
Pelecaniformes, Mesitornithiformes (highest first), and Piciformes.
Caprimulgiformes, Procellariiformes, Apterygiformes, Struthioniformes, and Psittaciformes
(lowest first, the same applies below) show the least disaprity residuals (\cref{subfig:models_order}).

\begin{figure}[htbp]
  \centering
  \begin{subfigure}{0.48\textwidth}
    \centering
    \includegraphics[width=\textwidth]{model_comparison_order}
    \caption{}
    \label{subfig:models_order}
  \end{subfigure}
  \hfill
  \begin{subfigure}{0.48\textwidth}
    \centering
    \includegraphics[width=\textwidth]{model_comparison_family}
    \caption{}
    \label{subfig:models_family}
  \end{subfigure}
  \caption{Comparison of models of the phenotypic disparity of taxa (orders or famalies) as a
  function of diversity. Showing scatter plot of all taxa, graphs of four models, names of
  the five taxa with the highest and the five with the lowest phenotypic disparity.}
  \label{fig:models_comparison}
\end{figure}

For families, the power law receives the lowest AIC (-1201.08), followed by the
Hill equation (-1199.76), without significant difference
($\Delta AIC = 1.32$). The models are respectively $f(x) = 1 - x^{-0.1429}$
and $f(x) = \frac{(x - 1)^{0.3723}}{6.0303 + (x - 1)^{0.3723}}$.
Mitrospingidae, Vangidae, Oreoicidae, Ptiliogonatidae, and Modulatricidae show the
highest disparity residuals, and Apodidae, Rhinocryptidae, Caprimulgidae, Hydrobatidae,
and Procellariidae show the lowest (\cref{subfig:models_family}).

\subsection{Phylogenetic signal}\label{subsec:r-phylogenetic-signal}

The visual morphospace exhibit a highly significant phylogenetic signal ($p = 0.001$, 999 permutations). 
It yields an overall $K_{mult}$ of 0.4150. 
The first phylogenetically aligned component (PAC1) alone exhibited extreme univariate Blomberg's K ($K_{uni} = 2.0775, p = 0.001$).
The first 23 PACs show a $K_{mult}$ of 1.0129, and the value crosses the threshold of 1.0 at PAC24.
In the permutation test, the maximum permuted $K_{mult}$ is 0.3627, and the maximum permuted $K_{uni}$ of PAC1 is 1.3290.

\begin{figure}[htbp]
    \centering
    \includegraphics[width=\textwidth]{trait_correlation_heatmap}
    \caption{Alignment between the CNN's embedding space with biological traits, represented by the proportion of explained variance. The horizontal axis represents the components of PACA (Phylogenetically Aligned Component Analysis), sorted by phylogenetic signal strength. The vertical axis represents bird traits in AVONET \parencite{tobias2022avonet}. Numerical traits are quantified by Pearson's $R^2$, and categorical traits are quantified by the $\eta^2$ based on Kruskal-Wallis H-statistics. Blank cells indicate statistical insignificance of the correlation ($p > 0.05$). Trait abbreviations are defined as same as \textcite{tobias2022avonet}.}
    \label{fig:paca-avonet}
\end{figure}

\begin{table}[htbp]
  \centering
  \caption{The PACA axis explained maximum variance of each AVONET trait. Trait abbreviations are defined as same as \textcite{tobias2022avonet}. }
  \label{tab:paca-avonet}
  \begin{tabular}{lll}
    \toprule         
    Trait & Max explained variances  &  PACA axes \\ 
    \midrule                   
    Beak.Length\_Culmen & 0.186940 &            1 \\ 
    Beak.Length\_Nares  & 0.131211 &            1 \\ 
    Beak.Width          & 0.114404 &            9 \\ 
    Beak.Depth          & 0.139089 &            9 \\ 
    Tarsus.Length       & 0.235985 &            1 \\ 
    Wing.Length         & 0.322378 &            1 \\ 
    Kipps.Distance      & 0.309666 &            1 \\ 
    Secondary1          & 0.231521 &            1 \\ 
    Hand-Wing.Index     & 0.197398 &            2 \\ 
    Tail.Length         & 0.131632 &            1 \\ 
    Mass                & 0.058428 &            1 \\ 
    Habitat             & 0.251840 &            1 \\ 
    Migration           & 0.109775 &            2 \\ 
    Trophic.Level       & 0.061948 &           10 \\ 
    Trophic.Niche       & 0.265584 &            1 \\ 
    Primary.Lifestyle   & 0.323673 &            1 \\ 
    \bottomrule                 
  \end{tabular}
\end{table}

As shown in \cref{fig:paca-avonet,tab:paca-avonet}, early (approximately first 10) PACA axes (highest phylogenetic signal) show the strongest explained variance (darkest red/orange) with most traits. Correlations decrease quickly as PACA dimension index increases, while do not go to zero entirely. The vast majority of traits peak at PACA dimension 1, which is the single axis of maximum phylogenetic signal. Mass, Trophic.Level, and Migration show low max explained variance.

\subsection{Disparity through time}\label{subsec:r-disparity-through-time}

The disparity-through-time (DTT) analysis reveals an early-burst pattern for the evolution of
avian visual morphospace. As shown in \cref{fig:dtt}, the empirical data of relative disparity
deviates extremely from the expectation of Brownian Motion null simulation.

Following the K-Pg mass extinction ($\sim 66$ Mya), the empirical disparity immediately increases
at an intense rate. This curve is consistently higher than the mean null expectation throughout
the Paleogene and Neogene, resulting in a positive Morphological Disparity Index (MDI).

Crucially, the empirical curve achieves approximately 50\% of the modern disparity shortly after
the K-Pg mass extinction, followed by a relative deceleration towards the present. Conversely,
the BM model requires significantly more time to reach the same disparity.

\begin{figure}[htbp]
    \centering
    \includegraphics[width=0.65\textwidth]{dtt_null_test_combined}
    \caption{Relative disparity-through-time (DTT) plot for avian visual morphospace.
    The solid blue line represents the empirical disparity, while the solid grey line
    indicates the mean expectation under a multivariate Brownian Motion (BM) null model.
    Both empirical and simulational data are normalised relative to the disparity of
    the extant time (0 Mya). The K-Pg boundary (66 Mya) is marked as the red dash line.}
    \label{fig:dtt}
\end{figure}

